\section{Verification and open questions}\label{sec:verify}

Three steps of the proofs rest on exact computer calculations. Apart from
these and from decimal values of closed forms, the proofs are done by hand.
In the proof of Lemma~\ref{lem:fisher-moments}, each moment that we need lies
in a space of sequences in $N$ of dimension at most 10, so its closed form
follows from its values at $N=1,\dots,10$. The script
\texttt{verify\_sticky\_fisher.py} computes these moments exactly for
$N=1,\dots,12$, with $q$ as a symbol, and checks with sympy the identities
for $\Var X^2$ and $R^2_N$ that follow from them. The proof of
Proposition~\ref{prop:triangles} uses 2 characteristic polynomials and 3
resultants of such polynomials, which \texttt{verify\_examples.py} computes
with sympy. The
formula of Proposition~\ref{prop:thirdorder}(i) is proved by hand, and
\texttt{verify\_second\_order.py} evaluates it exactly on the arc of 3 states
and on the whole 4-cycle, which gives the values $\frac54\sqrt2-2$ and
$-\frac9{16}$ in Remark~\ref{rem:C4}.

All other computations are checks at finite size, in Python and in double
precision unless we say otherwise. None of the proofs use them, and the
statements that we call numerical rest on them alone. There is one script
for each part of the paper. The script \texttt{verify\_first\_order.py}
covers Sections~\ref{sec:model} and~\ref{sec:firstorder}. The script
\texttt{verify\_second\_order.py} covers Section~\ref{sec:secondorder},
Proposition~\ref{prop:K3start} and Appendices~\ref{app:torus}
and~\ref{app:laplace}, and \texttt{verify\_examples.py} covers the rest of
Section~\ref{sec:examples}. Sections~\ref{sec:sticky} and~\ref{sec:fisher}
and Appendices~\ref{app:hdp} and~\ref{app:fisher} are checked by
\texttt{verify\_sticky\_fisher.py}, and Section~\ref{sec:planar} and
Appendix~\ref{app:planar} by \texttt{verify\_planar.py}. Each script computes
the numbers in the paper in 2 independent ways where it can, and compares
both with the printed values. For example, the run expansion of
Lemma~\ref{lem:runs} agrees with a list of all words and with a dynamic
program of the chain in exact rational arithmetic, and the torus
representation behind Theorem~\ref{thm:local} agrees with the run sum to
$7\times10^{-13}$. On faces with 2, 3 and 4 states of a random generator the
local law of Theorem~\ref{thm:local} is accurate to $1\%$ by $T=40$, and the
ratio in Lemma~\ref{lem:discrete} is $1-2.7\times10^{-5}$ at $T=10$ and
$N=10^6$. The constants $K_F$ are computed from the tilted generator in
double precision and from the Hessian of the Perron root in 30-digit
arithmetic. We located the crossings $M_F=M_G$ on the 4-cycle from the exact
face maxima up to $N=10^{14}$, and they agree with
Theorem~\ref{thm:crossing} and with the next term in
Proposition~\ref{prop:thirdorder}. For a vertex against an arc of 3 states at
$N=10^{14}$, the remainder in the crossing equation times $T_N$ is
$-0.2333$, while $c_F-c_G=\frac54\sqrt2-2=-0.2322$. The window of
Proposition~\ref{prop:K3start} was computed exactly up to $N=10^{20}$, and its
left end agrees to $3\times10^{-5}$ with the second-order equation. For the
directed 3-cycle of Proposition~\ref{prop:threecycle} the script follows the
first crossing up to $N=10^{12}$ and the second up to $N=10^8$, and an exact
dynamic program at $N=300$ shows the 3 successive modes.

On the graphs of Section~\ref{sec:examples} the exact law of the composition
comes from a dynamic program of the whole chain, which agrees with an
enumeration of all $d^N$ words in exact arithmetic to about $10^{-15}$, and
the face maxima come from the run expansion. For the planar walk the
three-level form of Lemma~\ref{lem:threelevel}, Fourier inversion, a dynamic
program and an enumeration of all $4^N$ words agree to $2\times10^{-13}$ for
$N\le10$, and the crossings up to $N=10240$ come from the three-level form.
The sticky dynamic program agrees with an enumeration of all paths to
$10^{-16}$, and the probabilities of the direct jump come from Monte Carlo
with $10^6$ draws and from quadrature in 2 coordinate systems. The Fisher
information is computed from the run formula of Paper~A and from its
recurrence with a complex step, and the 2 agree with exact rational
arithmetic for $N\le14$ and with 30-digit arithmetic for $N\le400$. The
script \texttt{verify\_all.py} runs all 5 scripts in a few minutes, and with
the option \texttt{--full} it also runs the largest cells. The scripts are in
the folder \texttt{paperC} of the Problem 131 folder, and their output is in
\texttt{paperC/data}. Before each new computation we wrote a prediction, and
a condition that would refute it, in the ledgers \texttt{ledger\_C1.md} to
\texttt{ledger\_C5.md} in \texttt{paperC/data}. The outcome is recorded below
each prediction, also when it refuted the prediction. None of the refuted
predictions concerns a statement that we call a theorem.

Some of the finite and algebraic steps are also proved in Lean 4, in the
folder \texttt{lean/PaperC} of the Problem 131 folder. For
Sections~\ref{sec:model} and~\ref{sec:firstorder} these are the run
expansion of Lemma~\ref{lem:runs}, the zero set in Example~\ref{ex:star}, the
lower bound and the value at $\alpha^*_F$ in Proposition~\ref{prop:rate}(a),
the form in part (c) for symmetric $Q$ at interior points, and the finite
statements behind Theorem~\ref{thm:hull}(b) and
Lemma~\ref{lem:symmetric}(c). For Section~\ref{sec:secondorder} they are the
identity of Proposition~\ref{prop:pair}, the exact parts of
Proposition~\ref{prop:threecycle}(i) and~(ii), and the offsets $b_{FG}$ in
Corollaries~\ref{cor:paperA} and~\ref{cor:paperB}. For
Section~\ref{sec:examples} they are the exit rates of the faces of $K_d$, of
the arcs of cycles and of the end-arcs of paths,
Proposition~\ref{prop:complete}(b)--(d), the window of
Proposition~\ref{prop:K3start}(a) and~(b), Lemma~\ref{lem:arcs}(b) and~(c),
the cascade of Theorem~\ref{thm:cycle}(b) and~(c) among the arcs and the
whole cycle, the constants $K_{\rm pair}$ and $K_{\rm arc3}$ of
Remark~\ref{rem:C4}, parts of Proposition~\ref{prop:triangles}, and the
winners on the paw in Remark~\ref{rem:nonnested}. For
Sections~\ref{sec:sticky} to~\ref{sec:fisher} they are
Propositions~\ref{prop:thresholds}(b)--(d) and~\ref{prop:directjump}(a)
and~(b), the finite bounds of Proposition~\ref{prop:mass}, the identity
$o^{(0)}_N=2S_N(u)c_N$ of Proposition~\ref{prop:decomp}, and
Proposition~\ref{prop:eta}. The header of each Lean file lists what it
proves. The local laws, the crossings and the other asymptotic results are not
in Lean.

Several questions remain open. Theorem~\ref{thm:local} holds on compact
subsets of the open face. When some $T\alpha_i$ is $O(1)$ there is a
boundary layer, and the rare-coordinate crossover of Paper~B is the case of
the complete graph. A local law that holds uniformly up to the boundary of
a face is open (Example~\ref{ex:star}). For a reducible face,
Theorem~\ref{thm:first}(ii) only places the power of $T$ in $M_F$ between
$k_F-k^*$ and $k_F-1$. Heuristically it is $(k^*-1)/2+k_F-k^*$, and
Proposition~\ref{prop:threecycle} shows that the coefficient of
$\log\log N$ in a crossing then changes. We have no formula for $c_F$ in
terms of trees or cycles of the graph, and for non-reversible faces we have
no tree formula for $K_F$ like Proposition~\ref{prop:constant}(b).
Theorem~\ref{thm:crossing} puts all zeros in an interval of length
$O(1/L)$ but does not show that there is only one. Papers~A and~B and
Proposition~\ref{prop:decomp} prove uniqueness with polynomial arguments
that do not extend to general faces. For linear functions of the
composition, such as the other endpoints of the planar walk, face selection
is only heuristic here. At finite $N$ the order of the phases can differ
from first order when a face that never wins lies close to the hull, as on
the path with 4 states (Remark~\ref{rem:paths}). We do not know a condition
on the hull that rules this out. For the planar walk a heuristic expansion
suggests that the next term of Theorem~\ref{thm:fourdir} comes from
replacing $T$ by $T+\frac1{4r}-\frac1{2(r+s)}$, and the crossover near
$s=2r$, where $\lvert s-2r\rvert\log N=O(1)$, is open. For sticky priors we
do not treat exit probabilities drawn at random as in~\cite{fox2011}, or the
probability of the direct jump for $d\ge4$, where triples of states enter
the criterion. Finally, we do not compute the constant in
Theorem~\ref{thm:fisher}, and at fixed $q$ we prove only a lower bound for
$2(N-1)pq\,e_N(q)$.
